\newgeometry{left=1.5cm,right=1.5cm,top=2cm,bottom=2.5cm}
\fancyfootoffset{0pt}% <- recalculate \headwidth

\part*{PPD2, Gruppenarbeit zu Digitalität\\\normalsize
Severin Henzi, Simon Kronenberg, Malte Scheurer, Cyril Wendl}

Aufgrund der unterschiedlichen Unterrichtsfächer (Physik, Geografie, Geschichte und Informatik) liegen vier Unterrichtseinheiten vor, die sich zum Teil deutlich voneinander unterscheiden. Die Unterrichtseinheiten aus GG, GS und PH setzen beispielsweise stärker auf gemeinschaftliche Arbeitsaufträge, während die Unterrichtseinheit IN auch einen grossen Anteil Einzelarbeit umfasst. Entsprechend sprechen in Bezug auf die Gruppengrösse (Kriterium 3) mehrere Personen zu mehreren Personen (Mehr-zu-Mehr-Kommunikation) beziehungsweise bei der IN-Einheit nur im gemeinsamen Teil als Eins-zu-Mehr-Kommunikation.

In Bezug auf die Gelingensbedingungen lässt sich insgesamt feststellen, dass eine hohe Selbstständigkeit beim Einsatz der eigenen Geräte durch die SuS vorausgesetzt wird. Bei den Unterrichtseinheiten PH sind dabei explizite Programmierkenntnisse gefordert, während es bei den Unterrichtseinheiten IN, GE und GS eher um den Einsatz von Hilfs-Tools geht. Diese hohe Selbstständigkeit darf aus unserer Sicht auch erwartet werden, da dies fachübergreifend wichtige zu erlernende Kompetenzen sind. Nichtsdestotrotz ist es dabei wichtig, dass die SuS zur Orientierung in den Lernaufgaben gewisse Vorgaben, Anhaltspunkte oder vorgegebene Strukturen auffinden, die den Übungen ihren Rahmen geben. Sei es beim Arbeitsprozess selbst durch konkrete Anweisungen oder ungefähren Zeitangaben des Aufwandes oder aber durch eine gewisse Vorstrukturierung der kollaborativen Dokumente oder auch mittels der nötigen technischen Einstellungen, die dem vorgesehenen Lern-Setting entsprechen und dieses zielgerichtet unterstützen. Dies muss durch die Lehrperson im Vorfeld der Unterrichtseinheit sauber geplant und umgesetzt werden, da es insbesondere beim Unterricht mit digitalen Medien wesentlich zu einem guten Gelingen beiträgt. Der Einsatz von digitalen Aufträgen ersetzt die Präsenz der Lehrperson nicht, diese bleibt eine wichtige Unterstützung im Unterricht.

Für das Gelingen der Einheiten wird bei allen vier Unterrichtseinheiten ein gutes Klassenklima und eine Klassenkultur, in der gut zusammengearbeitet werden kann, vorausgesetzt. Insbesondere über die kollaborativen Elemente in den Aufgaben wird das Arbeiten im Team weiter gefördert und diese soziale Ressource weiter aufgebaut, was sich langfristig auszahlt, da ein gutes Lernklima einen wesentlichen Bestandteil des erfolgreichen Lernens von Jugendlichen ist. Speziell bei den Übungen mit Gruppenarbeiten lernen die SuS sich in der Gruppe gemeinsam auf die wichtigsten Inhalte zu einigen und diese entsprechend zu präsentieren, was ein wichtiger Teil von Kommunikation darstellt. 

Alle vier Unterrichtseinheiten wurden nach dem AVIVA-Prinzip konzipiert, beachten in ihrer Lernzielorientierung unterschiedliche, kognitive Stufen von Wissen und nutzen digitale Tools für technologiegestützte Lernprozesse. Bei den Unterrichtseinheiten GG und GS steht aus Digitalitäts-Sicht das gemeinsame Arbeiten in gleichen Dokumenten im Zentrum. Es wird eingesetzt, um in der Recherche-Phase Inhalte zu sammeln (sog. divergenter Prozess; Kriterium 7), welche danach in Form eines Mindmaps bzw. eines Padlets verdichtet werden (sog. Konvergenter Prozess; Kriterium 7). Die dabei eingesetzten Tools erlauben eine gewisse Strukturierung der Inhalte, die entweder durch die Lehrperson bereits vorgegeben wird oder durch die SuS selbständig gewählt werden kann. 

Damit der Unterricht reibungslos verläuft und kollaboratives Arbeiten in digitalen Umgebungen ermöglicht wird, ist eine funktionierende technische Infrastruktur eine grundlegende Voraussetzung. Dies umfasst sowohl die Bereitstellung geeigneter Geräte im Rahmen eines BYOD-Konzepts, als auch den Zugang zu relevanter Software und digitalen Plattformen, auf denen Unterrichtsmaterialien zentral hinterlegt und jederzeit abgerufen werden können. 
Die Lehrperson ist dabei auch technische Ansprechperson für die SuS und sollte über die entsprechenden Kenntnisse der eingesetzten Tools und Programme verfügen, um einen produktiven und störungsfreien Unterricht zu gewährleisten. 

Bei allen vier Unterrichtseinheiten lässt sich dabei eine divergente Phase und eine konvergente Phase erkennen (Kriterium 7). Die divergente Phase umfasst das Sammeln und Erarbeiten von Informationen, was vorliegend mittels Padlet (GG), Mindmap/Sutori (GS) oder im Falle von IN im Einzelauftrag beim Lösen der Aufgabe geschieht. Die konvergente Phase wird bei allen Unterrichtseinheiten im Sinne der Verdichtung und dem Zusammenführen der erarbeiteten Aufgaben und Inhalte im Plenum eingesetzt und findet synchron (Kriterium 1) und ortsgebunden (Kriterium 2) statt. Dies ist in vielerlei Hinsicht eine äusserst wichtige Phase, da einerseits der Abgleich mit der Lösung stattfindet, das erworbene Wissen präsentiert und plausibiliert wird und allfällige Missverständnisse, Fehler oder Fragen besprochen werden können. Digitale Tools (wie Padlet, Video oder Sutori) werden hier unterstützend eingesetzt und entfalten ihre Wirkung voll.  
Die vier Unterrichtsentwürfe zeigen verschiedene Ansätze für modernen, digital gestützten Unterricht. Während sich Informatik und Physik stärker auf algorithmische und experimentelle Methoden konzentrieren, arbeiten Geografie und Geschichte mit reflektierenden Methoden. Trotz unterschiedlicher Fächer gibt es deutlich Übereinstimmungen im didaktischen Konzept.
